% ============================================================
% FRANÇAIS — PRÉSENT DE L'INDICATIF
% Version graphique LaTeX/TikZ
% Compilation : XeLaTeX
% ============================================================
\documentclass[a4paper,10pt]{article}

\usepackage[a4paper,top=9mm,bottom=8mm,left=11mm,right=11mm]{geometry}
\usepackage{fontspec}
\usepackage[french]{babel}
\usepackage{microtype}
\usepackage{tikz}
\usetikzlibrary{arrows.meta,calc,positioning}
\usepackage[most]{tcolorbox}
\usepackage{tabularx}
\usepackage{array}
\usepackage{xcolor}
\usepackage[normalem]{ulem}

\setmainfont{Noto Sans}
\setsansfont{Noto Sans}

% ------------------------------------------------------------
% PALETTE — PRÉSENT = VERT
% ------------------------------------------------------------
\definecolor{Green}{HTML}{319B58}
\definecolor{GreenDark}{HTML}{17613A}
\definecolor{GreenLight}{HTML}{EAF6EC}
\definecolor{GreenLine}{HTML}{8FC9A3}
\definecolor{Blue}{HTML}{17458F}
\definecolor{BlueLight}{HTML}{E9F3FA}
\definecolor{Red}{HTML}{D52F35}
\definecolor{RedLight}{HTML}{FCEBED}
\definecolor{Gold}{HTML}{B27A00}
\definecolor{GoldLight}{HTML}{FFF3D7}
\definecolor{Ink}{HTML}{172A32}
\definecolor{Grey}{HTML}{64757D}
\definecolor{Border}{HTML}{B9D3C3}
\definecolor{White}{HTML}{FFFFFF}

\pagestyle{empty}
\setlength{\parindent}{0pt}
\setlength{\parskip}{0pt}
\setlength{\tabcolsep}{3pt}

% ------------------------------------------------------------
% CADRE EXTÉRIEUR
% ------------------------------------------------------------
\newcommand{\PageFrame}{%
\begin{tikzpicture}[remember picture,overlay]
  \coordinate (A) at ([xshift=8mm,yshift=-8mm]current page.north west);
  \coordinate (B) at ([xshift=-8mm,yshift=-8mm]current page.north east);
  \coordinate (C) at ([xshift=-8mm,yshift=8mm]current page.south east);
  \coordinate (D) at ([xshift=8mm,yshift=8mm]current page.south west);
  \def\r{1.7mm}

  \draw[Border,line width=.65pt]
    (A)
    -- ([xshift=-\r]B)
    arc[start angle=90,end angle=0,radius=\r]
    -- ([yshift=\r]C)
    arc[start angle=0,end angle=-90,radius=\r]
    -- (D) -- cycle;

  \draw[Green,line width=3.2mm] (A) -- (D);
\end{tikzpicture}%
}

% ------------------------------------------------------------
% BOÎTES
% ------------------------------------------------------------
\tcbset{
  enhanced,
  boxrule=.6pt,
  arc=2.8mm,
  left=4mm,right=4mm,top=2mm,bottom=2mm,
  coltext=Ink
}

\newtcolorbox{timelinebox}{colback=White,colframe=GreenLine}
\newtcolorbox{greenbox}{colback=GreenLight,colframe=GreenLine}
\newtcolorbox{bluebox}{colback=BlueLight,colframe=Border}
\newtcolorbox{redbox}{colback=RedLight,colframe=Border}
\newtcolorbox{whitebox}{colback=White,colframe=Border}
\newtcolorbox{goldbox}{colback=GoldLight,colframe=GoldLight,boxrule=0pt,arc=2.5mm,left=2.7mm,right=2.7mm,top=1.25mm,bottom=1.25mm}

% ------------------------------------------------------------
% PETITES ÉTIQUETTES
% ------------------------------------------------------------
\newcommand{\pill}[2][GreenLight]{%
  \tikz[baseline=(X.base)]{\node[fill=#1,rounded corners=2pt,inner xsep=4.2pt,inner ysep=1.6pt,font=\small\bfseries,text=Ink] (X) {#2};}%
}
\newcommand{\sectiontitle}[2]{{\large\bfseries\textcolor{#1}{#2}}}
\newcommand{\signalpill}[1]{%
  \tikz[baseline=(X.base)]{\node[fill=GoldLight,rounded corners=2pt,inner xsep=3.7pt,inner ysep=1.1pt,font=\scriptsize\bfseries,text=Ink] (X) {#1};}%
}

% ------------------------------------------------------------
% LIGNE DU TEMPS — PRÉSENT FRANÇAIS
% ------------------------------------------------------------
\newcommand{\Timeline}{%
\begin{tikzpicture}[x=1cm,y=1cm]
  \draw[-{Latex[length=3mm]},very thick,Ink] (0,0) -- (15.3,0);

  % NOW : flèche séparée pointant vers la ligne
  \draw[-{Latex[length=3mm,width=2mm]},very thick,Green]
    (9.3,1.10) -- (9.3,.10);
  \node[font=\bfseries\small,text=Green] at (9.3,1.32) {MAINTENANT};
  \fill[Green] (9.3,0) circle (3.2pt);

  % zone du présent : autour du moment présent
  \draw[Green,line width=3.2pt] (8.35,0) -- (10.25,0);
  \fill[Green] (8.35,0) circle (2.1pt);
  \fill[Green] (10.25,0) circle (2.1pt);

  \node[font=\scriptsize\bfseries,text=Grey] at (2.2,-.48) {PASSÉ};
  \node[font=\scriptsize\bfseries,text=Green] at (9.3,-.48) {PRÉSENT};
  \node[font=\scriptsize\bfseries,text=Grey] at (13.3,-.48) {FUTUR};

  \node[fill=GreenLight,rounded corners=2pt,inner xsep=5pt,inner ysep=2pt,font=\scriptsize\bfseries,text=GreenDark]
    at (4.25,.72) {action / état au présent};
\end{tikzpicture}%
}

\begin{document}
\PageFrame

% ------------------------------------------------------------
% TITRE
% ------------------------------------------------------------
\vspace*{1mm}
\begin{center}
  {\fontsize{25}{27}\selectfont\bfseries\textcolor{Ink}{PRÉSENT}}

  \vspace{0.75mm}
  {\large\textcolor{Grey}{Le présent de l'indicatif}}

  \vspace{3mm}

  \begin{tcolorbox}[
    width=.88\textwidth,
    colback=GreenLight,
    colframe=GreenLight,
    boxrule=0pt,
    arc=3mm,
    halign=center,
    top=2.2mm,bottom=2.2mm
  ]
    {\large\bfseries\textcolor{GreenDark}{RADICAL + TERMINAISON DU PRÉSENT}}
    \quad
    {\small\textit{ex. : parl- + e, es, e, ons, ez, ent}}
  \end{tcolorbox}
\end{center}

\vspace{1mm}

% ------------------------------------------------------------
% TIMELINE
% ------------------------------------------------------------
\begin{timelinebox}
  \sectiontitle{Green}{La ligne du temps}

  \vspace{0.75mm}
  \begin{center}
    \Timeline
  \end{center}

  \vspace{-1mm}
  \begin{center}
    \small
    Le \textbf{présent} situe un fait \textbf{au moment présent},
    ou le présente comme \textbf{habituel, général ou durable}.
  \end{center}
\end{timelinebox}

\vspace{1.35mm}

% ------------------------------------------------------------
% FORMATION
% ------------------------------------------------------------
\begin{greenbox}
  \sectiontitle{Green}{Formation — exemple : parler}

  \vspace{0.75mm}
  \begin{center}
    \begin{tabular}{@{}c@{\qquad}c@{\qquad}c@{\qquad}c@{\qquad}c@{\qquad}c@{}}
      je & tu & il / elle / on & nous & vous & ils / elles\\[-.2mm]
      \textcolor{Green}{parle} & \textcolor{Green}{parles} & \textcolor{Green}{parle} &
      \textcolor{Green}{parlons} & \textcolor{Green}{parlez} & \textcolor{Green}{parlent}
    \end{tabular}
  \end{center}

  \vspace{0.5mm}
  \begin{center}
    \small\textcolor{Grey}{Pour les verbes en \textbf{-er}, les terminaisons sont :}
    \quad
    \textbf{-e \; -es \; -e \; -ons \; -ez \; -ent}
  \end{center}
\end{greenbox}

\vspace{1.35mm}

% ------------------------------------------------------------
% USES
% ------------------------------------------------------------
\begin{greenbox}
  \sectiontitle{Green}{Quand utilise-t-on le présent ?}

  \vspace{0.75mm}
  \begin{tabularx}{\textwidth}{@{}p{.31\textwidth} X@{}}
    \pill{ACTION EN COURS}
      & Action qui se déroule maintenant.\newline
        \textit{Je lis un livre.}\\[0.7mm]
    \pill{HABITUDE / RÉPÉTITION}
      & Action régulière ou habituelle.\newline
        \textit{Nous allons au marché le samedi.}\\[0.7mm]
    \pill{VÉRITÉ GÉNÉRALE}
      & Fait considéré comme généralement vrai.\newline
        \textit{L'eau bout à 100\,°C.}\\[0.7mm]
    \pill{ÉTAT / SITUATION DURABLE}
      & Situation, propriété ou état relativement stable.\newline
        \textit{Paul habite à Albi.}
  \end{tabularx}
\end{greenbox}

\vspace{1.35mm}

% ------------------------------------------------------------
% EXAMPLES
% ------------------------------------------------------------
\begin{bluebox}
  \sectiontitle{Blue}{Exemples}

  \vspace{0.75mm}
  \begin{tabularx}{\textwidth}{@{}X@{\hspace{7mm}}X@{}}
    \textbf{Je travaille aujourd'hui.}
    & \textbf{Nous mangeons souvent ensemble.}\\[-.5mm]
    \textit{Action en cours.}
    & \textit{Habitude.}\\[0.7mm]
    \textbf{La Terre tourne autour du Soleil.}
    & \textbf{Elle habite dans le Tarn.}\\[-.5mm]
    \textit{Vérité générale.}
    & \textit{Situation durable.}
  \end{tabularx}
\end{bluebox}

\vspace{1.35mm}

% ------------------------------------------------------------
% POINT DE VIGILANCE
% ------------------------------------------------------------
\begin{redbox}
  \textcolor{Red}{\bfseries !\quad À retenir :}
  en français, le présent peut traduire plusieurs temps anglais.\quad
  \textbf{Je travaille} peut signifier \textit{I work} ou \textit{I am working},
  selon le contexte.
\end{redbox}

\vspace{1.35mm}

% ------------------------------------------------------------
% MOTS REPÈRES
% ------------------------------------------------------------
\begin{goldbox}
  \begin{center}
    {\normalsize\bfseries\textcolor{Gold}{MOTS REPÈRES}}
    \hspace{3mm}
    \signalpill{maintenant}\hspace{1mm}
    \signalpill{aujourd'hui}\hspace{1mm}
    \signalpill{souvent}\hspace{1mm}
    \signalpill{toujours}\hspace{1mm}
    \signalpill{habituellement}

    \vspace{0.75mm}

    \signalpill{chaque jour}\hspace{1mm}
    \signalpill{tous les jours}\hspace{1mm}
    \signalpill{le lundi}\hspace{1mm}
    \signalpill{en général}
  \end{center}
\end{goldbox}

\end{document}
